% ============================================================
%  第6章 OpenCL主机端内存模型
%  源文件：content/chapter6/
% ============================================================

为了保证在各种硬件上都具有良好的可移植性，OpenCL提供了一个较为完善的抽象内存模型。这个内存模型足以应对各种硬件设备，提供足够强的内存序保证，以表达对开发者重要的各类并行，并且能够并行的完成。抽象的内存模型也作为连接编程者和硬件的重要接口。编程者可以基于模型规则进行代码的书写，而无需考虑设备端在执行内核时将如何处理这段内存。硬件供应商在实现其SDK运行时部分时，也需要将其硬件的具体组件映射到内存模型中。并且提前定义组件映射关系，即可保证编程者只能和相应的组件进行互动。

之前的章节中，已经涉及到一些关于OpenCL内存模型的内容。例如，在第3章和第4章已经使用到的数组和图像内存。也介绍了内存区域，例如：全局和局部。本章和第7章将要讨论内存模型更多的细节。我们将内存模型分为两个部分：主机端内存模型和设备端内存模型。设备端内存模型将在下章进行讨论，其内容与“内核在运行时如何使用内存对象和其他数据”相关。

OpenCL设备包括GPU和其他加速器，这些设备上的内存系统与CPU所使用的主存是分开的。通常，OpenCL主机端内存模型的一致性是松散的，这使得全局内存同步只能使用事件来完成。OpenCL 2.0标准中添加了对内存一致性的保证，其借鉴了C/C++11和Java类似的机制。

为了在系统中支持多个离散内存系统，以及各种一致性模型，这就使得OpenCL内存对象的定义与主机CPU的内存有所不同。需要使用对应API来完成，数据从CPU指针搬入OpenCL内存对象，或是CPU指针接收OpenCL内存对象的数据的操作。对于OpenCL内存对象来说，对其尤为重要的是上下文对象，而非设备对象。所以，数据在转移的时不会专门去指定具体的设备。这部分的工作就交由运行时来完成，运行时需要确保正确的数据，在正确的时间内出现在正确的位置上。

本章我们先要聊一下OpenCL中定义的内存对象的类型，然后在了解一下如何使用主机端API来管理内存对象。


% ---- 6.1-chinese.md ----
\section{内存对象}

OpenCL定义了三种内存对象——数组，图像和管道——这几种内存对象可以通过主机端的API进行创建。数组和图像内存对象上存储的数据，可以在主机端和设备端进行随机访问。管道对象上的数据对象只能在内核端先进先出(FIFO)，并且主机端无法访问这些数据。

数组对象可以看做为CPU上的一维数组，并且其分配过程与C的malloc()函数类似。数组对象可以包含任何标量数据，向量数据或自定义结构体。数据在数组中是顺序存储的，这样OpenCL内核就能以随机访问的方式对数组进行访问(就如同C中的一维数组)。

图像对象就有些不同，其数据的布局或存放方式在硬件上进行过一些优化，这样指针就很难直接一个一个的访问对应的数据，并且硬件上的数据布局方式对于开发者来说是不可见的。这样，内核端只能使用内置函数对图像对象进行访问。因为GPU设计之初就是为了处理图形任务，所以GPU对图像数据访问效率已有较高优化。图像有三个优势：

\begin{enumerate}
\item GPU上的层级缓存和数据流结构就是为了优化访问图像类型数组所准备
\item GPU驱动会在硬件层面上优化图像数据的排布，从而提升访问图像数据的效率，尤其是二维图像模式
\item 硬件支持图像是一个很复杂的数据访问过程，在这个过程中硬件会将一些存储的数据进行压缩
\end{enumerate}

下面的几个子节，将分别对这三种内存对象进行更为详尽的描述。

\subsection{数组对象}

数组对象的分配有点类似使用malloc()，其分配方式非常简单。创建数组对象只需要提供上下文对象，数组大小，以及一些标识就可以。创建数组对象的API为\texttt{clCreateBuffer()}。

\begin{codebox}[cpp]
cl_mem
clCreateBuffer(
  cl_context context,
  cl_mem_flags flags,
  size_t size,
  void *host_ptr,
  cl_int *err)
\end{codebox}

函数会返回一个数组对象，如果需要将错误码传出，则需要传入最后一个参数。flags参数可以将数组配置成只读或只写的数据，以及设置其他分配选项。例如，下面的代码，我们就创建了一个只读的数组对象，其与主机端数组a使用相同的存储位置，二者也具有同样的大小。这里，我们将详细讨论一些分配选项在之后的章节(例如：CL\_\allowbreak{}MEM\_\allowbreak{}USE\_\allowbreak{}HOST\_\allowbreak{}PTR)。错误码将从err传出，对应的错误码在OpenCL标准文档中都有定义。通常OpenCL函数执行成功，都会以CL\_\allowbreak{}SUCCESS作为错误码返回。

\begin{codebox}[cpp]
cl_int err;
int a[16];

cl_mem newBuffer = clCreateBuffer(
  context,
  CL_MEM_READ_ONLY | CL_MEM_USE_HOST_PTR,
  16 * sizeof(int),
  a,
  &err);

if (err != CL_SUCCESS){
  // Handle error as necessary
}
\end{codebox}

OpenCL也支持\emph{子数组对象}(subbuffer)，也就是可以将单独的数组对象再进行划分为更小的数组对象，这些数组对象可以相互覆盖，可以读或写和拷贝，以及和其父数组以相同的方式使用。注意有覆盖和包含关系的子数组对象，会让其父数组对象结构变的更加复杂，并且在实际使用过程中这种情况会造成一些未定义的行为。

\subsection{图像对象}

OpenCL中，图像与数组对象有三点不同：

\begin{enumerate}
\item 图像数据排布的不透明性，使其不能直接在内核中使用指针进行数据读取
\item 多维结构
\item 图像对数据成员有一定的要求，并不能像数组那让接受任意数据类型
\end{enumerate}

图像对象之所以在OpenCL中存在，因为GPU硬件设计已经对图像的存储方式进行过优化，这样会让设备在访问图像数据时的效率更高。访问图像的内置函数，并不能像数组对象那样提供各种方式的访存模式，不过其能作用就是能让一些滤波方式在硬件方面得到很好的支持，从而具有较高的效率。滤波操作中，会有基于一组以特定方式排布的像素进行变换，高效的内存访问则会对滤波操作进行加速。这些操作需要多次读取的长指令队列，不过有硬件的支持其执行效率很高。

图像数据可在内核端，通过特定的函数进行访问(第7章设备端内存模型会详细讨论)。主机端访问图像对象的方式与数组对象没有太大的区别，不过对图像对象操作的主机函数支持多维度的寻址。\texttt{clEnqueueReadImage()}更像\texttt{clEnqueueReadBufferRect()}，而非\texttt{clEnqueuReadBuffer()}。

图像对象和数组对象最大的区别，就是图像对象支持的格式。图像格式包括通道序和通道类型。通道顺序定义了通道的数量及其排列次序——例如,CL\_\allowbreak{}RGB，CL\_\allowbreak{}R或CL\_\allowbreak{}ARGB。通道类型就是要选择通道内数据存储的格式，从CL\_\allowbreak{}FLOAT到更省存储空间的CL\_\allowbreak{}UNORM\_\allowbreak{}SHORT\_\allowbreak{}565(其将一个16比特字打包后放入内存)。当内核代码要访问这些数据时，读取到的数据结果都上转换成标准的OpenCL C类型数据。图像格式支持的列表可以通过\texttt{clGetSupportedImageFormats()}获取。

图像对象可通过\texttt{clCreateImage()}进行创建，其声明如下：

\begin{codebox}[cpp]
cl_mem
clCreateImage(
  cl_context context,
  cl_mem_flags flags,
  const cl_image_format *image_format,
  const cl_image_desc *image_desc,
  void *host_ptr,
  cl_int *errcode_ret)
\end{codebox}

context，flags和host\_\allowbreak{}ptr这些与创建数组对象所需要的参数一致。图像类型(image\_\allowbreak{}format)和图像描述符(image\_\allowbreak{}desc)参数定义了图像的维度，数据格式和数据分布。这种结构已经在第4章进行详细的描述(一个初始化图像对象的完整例子)。

\subsection{管道对象}

OpenCL 2.0支持一种新的内存对象——管道对象。管道对象的数据结构是FIFO结构，其用来将一个内核对象的数据传递给另一个内核对象。因为之前OpenCL标准中的内存模型十分松散，所以根本无法实现管道，因为没有办法在上一个内核结束前确定内存的状态。通道的意义就是为了在两个内核中共享一部分数据，并且保证这些共享数据的状态。这种标识对于处理器来说是一种趋势，对于支持管道对象的任意设备，至少有能力实现内核间共享数据的原子操作，并且必须要有一套内存一致模型来支持内存获取和释放语义。

可以设想一下，如果设备具有这样的能力，编程者都可以用一个数组对象来实现属于自己的“管道对象”。OpenCL 2.0的内存模型中这个方案是可行的，但这需要大量的编程工作。管道具有并发“生产者-消费者”机制，这种方式比起之前的标准，让很多问题变的简单了许多(例如，当每个工作项生成了大量的输出数据，就可以对这些数据进行打包处理)。在同一设备上执行“生产者-消费者”任务时，就可以使用管道对象，这样也允许硬件供应商能将这块内存映射到一个延迟较低的内存区域中。管道对象是不允许主机端对其进行读写的，所以访问管道对象属于设备端内存模型。

管道数据通常称为\emph{包}(packets)，其包含了OpenCL C或用户自定义的类型。创建管道对象的API为\texttt{clCreatePipe()}，其声明如下：

\begin{codebox}[cpp]
cl_mem
clCreatePipe(
  cl_context context,
  cl_mem_flags flags,
  cl_uint pipe_packet_size,
  cl_uint pipe_max_packets,
  const cl_pipe_properties *properties,
  cl_int *errcode_ret)
\end{codebox}

当创建一个管道，需要提供包大小(pipe\_\allowbreak{}packed\_\allowbreak{}size)和最大包数(pipe\_\allowbreak{}max\_\allowbreak{}packets)。如同其他创建内存对象的API，这个API也需要设置一些与内存相关标识。对于管道对象来说，只有设置CL\_\allowbreak{}MEM\_\allowbreak{}READ\_\allowbreak{}WRITE标识是合法的，其也是管道对象默认的标识参数。以后，管道对象不可在主机端访问，即便是编程者没有意识到，也需要使用CL\_\allowbreak{}MEM\_\allowbreak{}HOST\_\allowbreak{}NO\_\allowbreak{}ACCESS来隐式表明管道对象在主机端不可访问。



% ---- 6.2-chinese.md ----
\section{内存管理}

使用OpenCL主机端API创建内存对象，这些内存对象都在全局空间内分配，可以被上下文上中所有的设备使用。虽然，OpenCL中只设置了一块全局内存，不过实际使用的很多异构系统中具有很多设备，这些设备对共享地址空间有严格的要求，并且可能不同设备的内存在物理设备上是分开的——比如CPU的内存和离散GPU的内存。这样的话，在内核运行前，运行时可能就要将数据在各种设备的内存间进行拷贝。即使在共享内存系统中，数据不需要重复拷贝，不过内核所需要的数据还是要拷贝到设备的层级缓存中。当存在数据拷贝，那么拷贝阶段某一设备上的内存可能是不一致的，即它与上下文中的其他设备所看到的内存不一致。存在有潜在的拷贝和不一致状态，那么如何保证最后一次复制的数据就是我们期望的数据呢？

后面的章节中，我们将讨论使用细粒度\emph{共享虚拟内存}(SVM)的内存序和可见性。现在我们假设我们使用的内存对象是默认内存对象，即\emph{非共享虚拟内存}(non-SVM)。当使用默认内存对象时，OpenCL松散的一致性模型不会让多个内核对在同一时间，对同一个内存对象进行操作。没有其他内核进行数据的修改，就保证了直到该内核执行完成，其他内核才能可见已经修改的内存。如果没有这种机制，那么运行时为一个内存对象创建了多个副本，两个内核对里面的内容进行修改，也就等同于对同一内存进行修改，那么最后得到结果必然是其中一个内核的结果覆盖另一个的结果；如果我们在让另外一个内核再去修改这块内存，那么这种修改就是会导致未定义的行为。除了需要编程者对自己的程序负责，运行时也要保证正确的数据在正确的时间，出现在正确的位置。将松散的内存模型和运行时对内存管理的责任结合在一起，就能让执行更加可移植，并且对编程者的影响降到最低。

另外，OpenCL标准的设计者也知道在，数据传输的低效，并且数据搬运会降低应用的性能。因此，OpenCL提供了一些命令，允许编程者设置“如何”，以及“在哪里”开辟这段内存；以及“在哪里”和“在何时”进行数据搬运。根据运行OpenCL系统的差别，不同方式分配方式和搬运方式对于应用的性能有很大的影响。下面两节中，就来聊聊怎么分配内存对象，以及转移内存对象上的数据——基于数组对象(虽然我们使用数组对象来做例子，不过也可以对图像对象进行相同的操作)。管道对象和数组与图像对象不同，其数据不能在主机端进行直接访问，所以就不存在搬运的问题。

\subsection{管理普通内存对象}

回到内存对象的穿件，我们使用\texttt{clCreateBuffer()}，并传入一系列flags，以及host\_\allowbreak{}ptr。这里在展示一些\texttt{clCreateBuffer()}的声明：

\begin{codebox}[cpp]
cl_mem
clCreateBuffer(
  cl_context context,
  cl_mem_flags flags,
  size_t size,
  void *host_ptr,
  cl_int *errcode_ret)
\end{codebox}

通过flags传入一些选项，就能告诉运行时应该“怎么”和“在哪里”分配该数组空间，并且host\_\allowbreak{}ptr用来初始化数组对象，或将主机端的数据拷贝到内存对象中。本节将了解当没有与分配相关的选项传入flags时，内存对象时如何工作的。下节中将了解，编程者在设置这些分配选项时，内存如何在硬件上进行分配。

通常，OpenCL不会指定在某个物理存储空间内分配内存对象——其将默认内存都看做"全局内存"。例如，运行时会决定是在CPU主存上分配内存，或是在离散GPU的内存上分配内存。这很像为数据创建多个空间，并且根据需要进行转移最后的数据副本。

当创建一个对象时，可以让主机端提供一个合法指针host\_\allowbreak{}ptr用来初始化创建的内存对象，并且将CL\_\allowbreak{}MEM\_\allowbreak{}COPY\_\allowbreak{}HOST\_\allowbreak{}PTR作为实参传入flags，这就指定了内存对象中的数据是从主机端进行拷贝得到。创建内存对象的过程不会产生对应的事件对象，我们可以假设内存对象是在拷贝完host\_\allowbreak{}ptr上的数据后才返回的。图6.1中展示运行时如何在数组对象创建和初始化、传递内核参数，以及从内存对象中读回数据中，数据转移的过程。

\myGraphic{0.43}{images/chapter6/6-1.png}{图6.1 一个示例：缓冲区在主机端创建并初始化，在设备端参与计算，再传回主机端。注意，运行时也可以直接在设备端创建并初始化该缓冲区。(a)在主机内存中创建并初始化缓冲区；(b)内核执行前隐式地把数据从主机传到设备；(c)显式地把数据从设备拷回主机指针。}

一般来说，从一个设备到另一个设备的转移数据的效率很低——第8章中，我们将讨论基于现代CPU和GPU，比较不同数据互传方式的内存带宽。为了让编程者更高效的转移数据，OpenCL提供了一些API(使用不同方式)专门用来进行数据搬移。其中最佳的选择要依赖于所实现的算法，以及目标系统的特点。

第一组命令就是显式的从主机端或设备端，将数据拷贝到设备端或主机端，其运行方式如图6.2所示。这两个命令为\texttt{clEnqueueWriteBuffer()}和\texttt{clEnqueueReadBuffer()}。\texttt{clEnqueueWriteBuffer()}的声明如下：

\begin{codebox}[cpp]
cl_int
clEnqueueWriteBuffer(
  cl_command_queue command_queue,
  cl_mem buffer,
  cl_bool blocking_write,
  size_t offset,
  size_t size,
  const void *ptr,
  cl_uint num_events_in_wait_list,
  const cl_event *event_wait_list,
  cl_event *event)
\end{codebox}

\myGraphic{0.44}{images/chapter6/6-2.png}{图6.2 使用显式读写命令搬运数据。(a)在设备内存中创建一个未初始化的缓冲区；(b)内核执行前显式地把数据从主机传到设备；(c)执行完成后显式地把数据从设备传回主机。}

\texttt{clEnqueueWriteBuffer()}和\texttt{clEnqueueReadBuffer()}的声明很类似，除了使用blocking\_\allowbreak{}write和blocking\_\allowbreak{}read。声明中可以看出，数据是在buffer和ptr之间进行传输。其中写命令就将主机端的数据在设备端进行拷贝(在全局内存中进行备份)，并且读命令就将设备端的数据在主机端进行拷贝。注意这里需要有命令队列参与。这样就需要指定设备，对其内存副本进行初始化。当然，对于大多数应用的编程者，他们是知道数据要拷贝到哪个设备上去的，并且指定对应的命令让运行时知道应该向哪个设备进行数据拷贝，从而就能避免多余的拷贝过程。这种设计也让运行时能更快的将相应的数据传递到对应的设备，以便在内核执行时供内核使用。

编程者可以通过设置offset和size，来决定从哪里拷贝多少byte个数据。注意ptr是作为主机端读或写的起始地址，offset用来决定数组对象的起始地址，而ptr的初始地址则由编程者决定。

这种数据传输是可以是异步的。当使用异步方式调用\texttt{clEnqueueReadBuffer()}时，函数返回时我们无法知晓拷贝过程是否完成，直到我们通过同步机制——通过事件机制，或调用\texttt{clFinish()}。如果我们想在函数返回时就完成拷贝，则需要将CL\_\allowbreak{}TRUE作为实参传入blocking\_\allowbreak{}write或blocking\_\allowbreak{}read中。这样，下面的代码中A 和 B 不保证会打印出不同的结果(即使outputBuffer最终应该是这个结果)。对于C的打印是能确定从outputBuffer拷贝出来的结果。

\begin{codebox}[cpp]
int returnedArray[16];
cl_mem outputBuffer;
cl_event readEvent;

// Some code that fills the returned array with 0s and invokes kernels
// that generates a result in outputBuffer
printf("A: %d\n", returnedArray[3]);

clEnqueueReadBuffer(
  commandQueue,
  outputBuffer, /* buffer */
  CL_FALSE, /* nonblocking read*/
  0,
  sizeof(int) * 16,
  returnedArray, /* host ptr */
  0,
  0,
  &readEvent);

printf("B: %d\n", returnedArray[3]);
clWaitForEvents(1, &readEvnet);
printf("C: %d\n", returnedArray[3]);
\end{codebox}

同步对于OpenCL内存模型尤为重要。修改中的内存不保证可见，且不保证内存状态的一致性，直到用一个事件来表明该命令结束(我们将在后面章节中讨论SVM内存的不同)。主机指针和设备内存间的数据互传，我们不能在拷贝的同时，对内存数据进行其他的操作，直到我们确定互传完成。仔细阅读过OpenCL标准文档的同学可能会发现，与设备内存相关的是上下文对象，而非设备对象。通过\texttt{clEnqueueWriteBuffer()}入队一个命令，直到其完成，过程中不能确定数据是否已经完全搬运到设备上，而能确定的是，主机端的数据已经开始进行转移。

与其他API不同，数据互传命令也可以指定为是同步的。我们只需要简单的将之前的调用进行修改即可：

\begin{codebox}[cpp]
clEnqueueReadBuffer(
  commandQueue,
  outputBuffer,
  CL_TRUE, // blocking read
  0,
  sizeof(int) * 16，
  returnedArray,
  0,
  0,
  &readEvent);
\end{codebox}

这样的调用，API需要保证设备端的数据完全传回主机端后才进行返回，并且在返回之后，主机端才能对读回的数据进行操作。

OpenCL提供了另一种命令进行主机和设备间的数据转换——\texttt{clEnqueueMigrateMemObjects()}，用来从当前地址(无论主机还是设备)转移到指定设备上。例如，如果一个数组创建和初始化使用的是CL\_\allowbreak{}MEM\_\allowbreak{}COPY\_\allowbreak{}HOST\_\allowbreak{}PTR时，可以调用\texttt{clEnqueueMigrateMemObjects()}显式的将数据转移到对应设备上。如果一个系统中有多个异构设备，那么该API也能用来进行不同设备间的数据数据交互。注意设备间数据交互不是使用\texttt{clEnqueueReadBuffer()}和\texttt{clEnqueueWriteBuffer()}，只有设备和主机端的数据进行交互时，才会使用这两个API。\texttt{clEnqueueMigrateMemObjects()}的声明如下：

\begin{codebox}[cpp]
cl_int
clEnqueueMigrateMemObjects(
  cl_command_queue command_queue,
  cl_uint num_mem_objects,
  const cl_mem *mem_objects,
  cl_mem_migrationg_flags flags,
  cl_uint num_events_in_wait_list,
  const cl_event *event_wiat_list,
  cl_event *event)
\end{codebox}

与之前的数据传输命令不同，\texttt{clEnqueueMigrateMemObjects()}需要使用内存对象数组作为实参传入，其可以通过一条命令转移多个内存对象。如同所有\texttt{clEnqueue*}开头的函数一样，该函数也能产生事件对象，指定依赖关系。当事件对象的状态设置为CL\_\allowbreak{}COMPLETE时，代表着相对应的设备端内存，已经传递到参数中的command\_\allowbreak{}queue命令队列上了。

除了显式告诉运行时进行数据转移，该命令也可以进行更高效的隐式转移。当编程者入队该命令后，恰好设备端在执行的任务与该数据转移命令没有任何关系(例如，内核执行时不包括该数据转移的内存)，数据转移会在任务执行时进行，从而隐藏了传输的延迟。这种隐式传递的方式\texttt{clEnqueueReadBuffer()}和\texttt{clEnqueueWriteBuffer()}也适用。

对于转移的内存对象，标准提供了一个标识：CL\_\allowbreak{}MIGRATE\_\allowbreak{}MEM\_\allowbreak{}OBJECT\_\allowbreak{}HOST。这个标识以为这告诉运行时，数据需要传输到主机端。如果设置该标识，那么传入的命令队列对象，将被API忽略。

\subsection{管理带有分配选项的内存对象}

之前的章节中，我们讨论的都是主机和设备端的数据互相进行拷贝，或转移设备端的内存。那么本章，我们将来讨论一下，如何在主机端和设备端直接访问物理上分离的内存区域。

OpenCL为程序员提供两种直接访问物理上分离的内存地址的方式，这两种方式会在主机端产生映射，可供主机端直接使用。创建设备端内存时，将CL\_\allowbreak{}MEM\_\allowbreak{}ALLOCL\_\allowbreak{}HOST\_\allowbreak{}PTR作为实参传入，运行时就会为主机端分配一个映射指针，供主机端代码进行操作。而CL\_\allowbreak{}MEM\_\allowbreak{}USE\_\allowbreak{}HOST\_\allowbreak{}PTR则告诉运行时，直接使用host\_\allowbreak{}ptr作为设备对象的空间。两个选项代表着两种不同的分配方式，而这两个选项则是互斥的，不可用于同一个内存对象。注意这里的CL\_\allowbreak{}MEM\_\allowbreak{}ALLOCL\_\allowbreak{}HOST\_\allowbreak{}PTR分配“可供主机端访问的映射内存”，这个概念其实是模糊的，其可能是通过主存链接主处理器，也有可能是真的能将设备上的一段内存映射到主机端。

在分配内存对象时，CL\_\allowbreak{}MEM\_\allowbreak{}COPY\_\allowbreak{}HOST\_\allowbreak{}PTR可以和CL\_\allowbreak{}MEM\_\allowbreak{}ALLOC\_\allowbreak{}HOST\_\allowbreak{}PTR一起使用，这样能创建映射到主机端的内存的同时，也能用主机端内存对内存对象进行初始化。不过，CL\_\allowbreak{}MEM\_\allowbreak{}COPY\_\allowbreak{}HOST\_\allowbreak{}PTR和CL\_\allowbreak{}MEM\_\allowbreak{}USE\_\allowbreak{}HOST\_\allowbreak{}PTR是不能同时使用的，因为在传入host\_\allowbreak{}ptr时，host\_\allowbreak{}ptr是已经分配好的一段内存，CL\_\allowbreak{}MEM\_\allowbreak{}USE\_\allowbreak{}HOST\_\allowbreak{}PTR代表使用host\_\allowbreak{}ptr作为设备端内存空间，其无法对自己进行初始化。

那么有没有一种选项组合，能完成内存数据在CPU端分配，并且设备可以直接对主机端数据进行访问。答案是肯定的，“确实有”。实际上这种方式的访存存在于一些系统中。一个系统中具有一个CPU和一个离散GPU，这种情况下GPU就需要通过PCIe总线对主存数据进行访问。当设备能用这种方式直接访问主机端内存，那么这种数据通常被称为\emph{零拷贝}(zero-copy)数据。

虽然使用CL\_\allowbreak{}MEM\_\allowbreak{}USE\_\allowbreak{}HOST\_\allowbreak{}PTR或CL\_\allowbreak{}MEM\_\allowbreak{}ALLOC\_\allowbreak{}HOST\_\allowbreak{}PTR将在主机端创建零拷贝数据(但这并不是OpenCL标准显式要求的)。当然，将CPU端的数据拷贝到设备端，供内核执行使用是最保险的。不过，在这种情况下，离散 GPU 可以直接从自己的显存中访问该缓冲区。实际上，如果使用CL\_\allowbreak{}MEM\_\allowbreak{}ALLOC\_\allowbreak{}HOST\_\allowbreak{}PTR，在主机端上分配出的指针，仅为设备端内存的映射。OpenCL中对于CL\_\allowbreak{}MEM\_\allowbreak{}ALLOC\_\allowbreak{}HOST\_\allowbreak{}PTR的描述，是能创建可供主机端访问的内存，并不是在主机端又分配了一段内存出来。

共享内存的系统或是单CPU环境中，使用CL\_\allowbreak{}USE\_\allowbreak{}HOST\_\allowbreak{}PTR可能就没有那些无谓的初始化拷贝，从而为应用用带来更好的性能。例如，设想将CPU作为设备，在CPU上执行内核：如果不指定CL\_\allowbreak{}MEM\_\allowbreak{}USE\_\allowbreak{}HOST\_\allowbreak{}PTR，应用将会另外开辟一段空间用来存放数据的副本。同样的情况也会出现在共享内存处理器上，比如加速处理单元(APU)。如果CPU和GPU共享同一块内存，那么在APU上是不是也要使用CL\_\allowbreak{}MEM\_\allowbreak{}USE\_\allowbreak{}HOST\_\allowbreak{}PTR呢？这个问题很难回答，因为这设计到相关内容的优化，这里只能说具体问题具体分析了。

APU的例子中，OpenCL运行时或设备端驱动，使用指定的分配标志，可能会对实际设备进行访存优化(例如：缓存与非缓存)；或是有另外的性能考量，例如\emph{非一致内存访问}(UNMA, nonuniform memory access)。APU上使用OpenCL时，CL\_\allowbreak{}MEM\_\allowbreak{}USE\_\allowbreak{}HOST\_\allowbreak{}PTR可能会将缓冲区视为可缓存且完全一致(coherent)。这将导致GPU访存效率降低，因为GPU先要确定CPU缓存访问的优先级。APU上进行编程时，编程者需要了解不同的选项所创建出来的内存有何不同。

CL\_\allowbreak{}MEM\_\allowbreak{}USE\_\allowbreak{}HOST\_\allowbreak{}PTR和CL\_\allowbreak{}MEM\_\allowbreak{}ALLOC\_\allowbreak{}HOST\_\allowbreak{}PTR都能创建出主机端能够访问的数据，OpenCL也提供了一种映射机制，能够直接操作主机端或设备端数据，而无需显式使用API进行读取和写入。与\texttt{clEnqueueReadBuffer()}拷贝数据的方式不同，这里的映射实现并非意味着进行了拷贝。有了这种机制，零拷贝内存才算完美实现。设备端在修改零拷贝内存的过程对主机不可见，直到其完成修改才再对主机可见。

调用\texttt{clEnqueueMapBuffer()}即可对一个内存对象进行设备端的映射：

\begin{codebox}[cpp]
void *
clEnqueueMapBuffer(
  cl_command_queue command_queue,
  cl_mem buffer,
  cl_bool blocking_map,
  cl_map_flags map_flags,
  size_t offset,
  size_t size,
  cl_uint num_events_in_wait_list,
  const cl_event *event_wait_list,
  cl_event *event,
  cl_int *errcode_ret)
\end{codebox}

调用\texttt{clEnqueueMapBuffer()}时，其会返回一个可供主机端访问内存的指针。当\texttt{clEnqueueMapBuffer()}产生的事件对象的状态为CL\_\allowbreak{}COMPLETE时，意味着主机端可以安全使用返回指针，进行数据访问。与\texttt{clEnqueueWriteBuffer()}和\texttt{clEnqueueReadBuffer()}相同，\texttt{clEnqueueMapBuffer()}也可以设置同步方式，将CL\_\allowbreak{}TRUE作为实参传入blocking\_\allowbreak{}map时，函数将在主机端可以安全使用访存指针时，将产生的指针返回。

\texttt{clEnqueueMapBuffer()}还有一个map\_\allowbreak{}flags，这个标识可以设置的实参有：CL\_\allowbreak{}MAP\_\allowbreak{}READ，CL\_\allowbreak{}MAP\_\allowbreak{}WRITE和CL\_\allowbreak{}MAP\_\allowbreak{}WRITE\_\allowbreak{}INVALIDATE\_\allowbreak{}REGION。CL\_\allowbreak{}MAP\_\allowbreak{}READ，主机只能对这块映射内存进行读取；CL\_\allowbreak{}MAP\_\allowbreak{}WRITE和CL\_\allowbreak{}MAP\_\allowbreak{}WRITE\_\allowbreak{}INVALIDATE\_\allowbreak{}REGION都表示主机端只能对该指针内容进行修改。CL\_\allowbreak{}MAP\_\allowbreak{}WRITE\_\allowbreak{}INVALIDATE\_\allowbreak{}REGION是带有优化的选项，其会指定整个区域将会被修改或者忽略，并且运行时在其修改完之前不会对中间值进行映射。这种方式就无法保证数据的一致性状态，而运行时对于内存区域的潜在访问要快于CL\_\allowbreak{}MAP\_\allowbreak{}WRITE。

当主机端对映射数据修改完毕，其需要进行反映射API的调用。反映射时需要将制定内存对象和映射出的指针传入该API。\texttt{clEnqueueUnmapMemObject()}的声明如下：

\begin{codebox}[cpp]
cl_int
clEnqueueUnmapMemObject(
  cl_command_queue command_queue,
  cl_mem memobj,
  void *mapped_ptr,
  cl_uint num_events_in_wait_list,
  const cl_event *event_wait_list,
  cl_event *event)
\end{codebox}

反映射时需要将内存对象本身和映射指针传入该API中。如其他数据管理命令一样，当API产生的事件对象的状态为CL\_\allowbreak{}COMPLETE时，则代表数据更新完成。与其他API不同\texttt{clEnqueueUnmapMemObject()}没有阻塞参数。图6.3展示了一个内存对象从映射到反映射的过程。如果内存对象在设备端修改的时候进行映射，而设备却对其进行写入，这种情况会引发一些未定义行为；反之亦然。

\myGraphic{0.51}{images/chapter6/6-3.png}{图6.3 使用映射/反映射搬运数据。(a)在设备内存中创建一个未初始化的缓冲区；(b)把该缓冲区映射到主机地址空间；(c)解除该缓冲区在主机地址空间的映射。}

本节我们一直在讨论内存的标志，对应标志的实际行为，很大程度上还是依赖于指定硬件上的相关实现。为了给编程者在分配内存上的一些建议，我们将简要的描述一下AMD旗下的设备是如何处理这些标志的。默认的内存对象(并未提供任何标志)将直接在设备端分配。当提供CL\_\allowbreak{}MEM\_\allowbreak{}COPY\_\allowbreak{}HOST\_\allowbreak{}PTR或CL\_\allowbreak{}MEM\_\allowbreak{}USE\_\allowbreak{}HOST\_\allowbreak{}PTR标志时，如果设备端支持虚拟内存，那么这块内存将创建为主机端的页锁定内存(不可分页)，并且也能被设备作为零拷贝内存直接访问。如果设备不支持虚拟内存，那么就会像默认内存对象一样，在设备端创建内存。当编程者希望将数据分配在设备端，且能让主机端直接访问该内存的话，AMD提供了一种供应商指定扩展的标志——CL\_\allowbreak{}MEM\_\allowbreak{}USE\_\allowbreak{}PERSISTENT\_\allowbreak{}MEM\_\allowbreak{}AMD。当在AMD设备上使用了该标志，则可以在设备端分配内存，且主机端可以直接访问。



% ---- 6.3-chinese.md ----
\section{共享虚拟内存}

OpenCL 2.0的一个重要的修订就是支持共享虚拟内存(SVM)。共享虚拟内存属于全局内存，其相当于在主机内存区域上做的扩展，其允许上下文对象上的所有设备能和主机共享这块内存。SVM能将结构体指针作为参数传入内核，这相较之前的标准方便许多。例如，在未使用SVM之前，在主机端创建的链表结构无法传递到内核中，只能一个节点一个节点的传入。那么，2.0之后如何将链表实例通过传参的形式传入内核呢？将内存块拆分成小块的方式不适合OpenCL的内存模型，将其整块传入内核才是真正合理的方式。SVM将打破原先内存模型的限制。

OpenCL 2.0中定义了三种SVM类型：

\begin{enumerate}
\item 粗粒度数组SVM
\item 细粒度数组SVM
\item 系统细粒度SVM
\end{enumerate}

读者可以通过表6.1来了解，我们将要讨论的各种SVM。

\par\medskip\noindent
\begingroup\fontsize{9bp}{12.5bp}\selectfont\renewcommand{\arraystretch}{1.12}
\begin{center}
\begin{tabularx}{\textwidth}{XllX}
\toprule
 & 粗粒度SVM & 细粒度SVM & 系统细粒度SVM \\
\midrule
OpenCL对象 & 数组 & 数组 & 无(任意主机端类型) \\
共享粒度 & 数组 & 字节 & 字节 \\
分配API & clSVMAlloc & clSVMAlloc & malloc(或类似的C/C++函数或操作) \\
一致性 & 同步点 & 同步点和选择性原子操作 & 同步点和选择性原子操作 \\
显式同步设备端和主机端数据？ & 映射/反映射命令 & 无 & 无 \\
\bottomrule
\end{tabularx}
\end{center}
\endgroup
\medskip

粗粒度数组SVM可以与OpenCL内存对象进行虚拟地址共享。粗粒度SVM内存与非SVM内存的不同点在于，主机和设备可以共用虚拟内存指针。粗粒度SVM内存需要在主机端进行映射和反映射，以保证设备所做的最新更新对主机可见。为了完成这个功能，主机端线程需要调用\texttt{clEnqueueMapBuffer()}将指定的内存区域阻塞的进行映射。当映射完成后，内核就可以对该内存进行使用。当\texttt{clEnqueueMapBuffer()}返回时，内核对该内存的任何操作，对于主机都是可见的。

我们使用\texttt{clCreateBuffer()}创建非SVM内存，同样SVM也有其创建API——\texttt{clSVMAlloc()}，其声明如下：

\begin{codebox}[cpp]
void *
clSVMAlloc(
  cl_context context,
  cl_svm_mem_flags flags,
  size_t size,
  unsigned int alignmet)
\end{codebox}

与非SVM内存对象一样，SVM也可以通过标志指定为：只读，只写和可读写。alignment表示内存对象需要在该系统上以最少多少字节对齐。如果传入0，则因为这使用默认的对齐字节，那么将会是OpenCL运行时支持的最大数据类型的大小。与\texttt{clCreateBuffer()}返回cl\_\allowbreak{}mem不同，\texttt{clSVMAlloc()}返回的是一个void型指针。就像C函数malloc()一样，\texttt{clSVMAlloc()}也会返回一个非空的指针来表示内存分配成功，否则分配失败。

释放SVM内存需要使用\texttt{clSVMFree()}函数，其只需要传入对应的上下文对象和SVM指针即可。

\begin{codebox}[cpp]
void
clSVMFree(
  cl_context context,
  void *svm_pointer)
\end{codebox}

\texttt{clSVMFree()}函数的调用会瞬间结束，而不需要等待什么命令结束。将SVM内存使用\texttt{clSVMFree()}函数释放之后在进行访问，程序会出现段错误，这就和普通的C程序没有任何区别了。为了保证在一些列命令使用完SVM之后，再对SVM进行释放，OpenCL也提供了一种入队释放的方式：\texttt{clEnqueueSVMFree()}。

与粗粒度数组SVM不同，细粒度数组SVM支持的是字节级别的数据共享。当设备支持SVM原子操作时，细粒度数组SVM内存对象可以同时在主机端和设备端，对同一块内存空间进行读与写。细粒度数组SVM也可被同一或不同设备，在同一时间对相同的区域进行并发访问。SVM原子操作可以为内存提供同步点，从而能保证OpenCL内存模型的一致性。如果设备不支持SVM原子操作，主机端和设备端仍可读取同一区域，但写入必须发生在互不重叠的区域，否则会导致数据覆盖。

将CL\_\allowbreak{}MEM\_\allowbreak{}SVM\_\allowbreak{}FINE\_\allowbreak{}GRAIN\_\allowbreak{}BUFFER标志传入\texttt{clSVMAlloc()}，就能创建细粒度数组SVM对象。若要使用SVM原子操作，则需要将CL\_\allowbreak{}MEM\_\allowbreak{}SVM\_\allowbreak{}ATOMICS一并传入flags中。注意，CL\_\allowbreak{}MEM\_\allowbreak{}SVM\_\allowbreak{}ATOMICS只有在同时指定CL\_\allowbreak{}MEM\_\allowbreak{}SVM\_\allowbreak{}FINE\_\allowbreak{}GRAIN\_\allowbreak{}BUFFER时才有效，否则即为非法。

细粒度系统SVM是对细粒度数组SVM的扩展，其将SVM的范围扩展到主机端的整个内存区域中——开辟OpenCL内存或主机端内存只需要使用malloc()就可以。如果设备支持细粒度系统SVM，那么对于OpenCL程序来说，内存对象的概念就不需要了，并且内存传入内核将是一件很简单的事情(如同CUDA内核函数的调用一样)。

查看设备支持哪种SVM，可以将CL\_\allowbreak{}DEVICE\_\allowbreak{}SVM\_\allowbreak{}CAPABILITIES标识传入\texttt{clGetDeviceInfo()}中进行查询。OpenCL标准规定，如果支持2.0及以上标准，则至少要支持粗粒度数组SVM。



% ---- 6.4-chinese.md ----
\section{本章总结}

本章从主机端的角度讨论了OpenCL的内存模型。主机在内存模型中的角色，大多是分配和管理全局内存对象(数组、图像和管道)。我们也对内存分配标志做了详细的介绍，让编程者了解，如何分配内存和在哪分配内存；一些管理标志可以用来指导编程者，从哪里搬移数据，以及何时进行数据搬运。本章也对SVM内存进行了介绍，SVM的内容将会在下章讨论设备端内存模型时继续进行。


